\documentclass[12pt, a4paper]{report}
\input{common/volume-preamble}

\begin{document}
\ExplSyntaxOff
\renewcommand{\familydefault}{\rmdefault}

\chapter*{Logic Opening Cards Mockup}

\section*{Shared Logical Machinery}

\begin{modelbox}{Free \(\Sigma\)-Algebra and Evaluation}
\modeldepends{Signature; inductively generated set; unique readability}

\begin{modelcomponent}{Layer}
Shared bridge machinery. Parameter: a signature \(\Sigma\). This card is
instantiated by connective evaluation, term evaluation, parsing, and later by
initial-algebra presentations such as \(\mathbb N\).
\end{modelcomponent}

\begin{modelcomponent}{Signature and Algebra}
\[
  \Sigma\text{ is a set of operation symbols},
  \qquad
  \operatorname{ar}:\Sigma\to\mathbb N.
\]
A \(\Sigma\)-algebra is a carrier \(\lvert\mathcal A\rvert\) with operations
\[
  f^{\mathcal A}:\lvert\mathcal A\rvert^{\operatorname{ar}(f)}
  \to \lvert\mathcal A\rvert
  \qquad(f\in\Sigma).
\]
Seeds land in the carrier; homomorphisms are maps of algebras.
\end{modelcomponent}

\begin{modelcomponent}{Free Term Algebra}
\[
  T_{\Sigma}(X),
  \qquad
  f^{T}(t_1,\ldots,t_n):=f(t_1,\ldots,t_n),
  \qquad
  \eta:X\to T_{\Sigma}(X).
\]
\end{modelcomponent}

\begin{modelcomponent}{Homomorphism}
\[
  \operatorname{Hom}(h,\mathcal A,\mathcal B)
  \quad:\Longleftrightarrow\quad
  h(f^{\mathcal A}(\bar a))
  =
  f^{\mathcal B}(h(a_1),\ldots,h(a_n))
\]
for every \(f\in\Sigma\) and \(\bar a\in\lvert\mathcal A\rvert^n\).
\end{modelcomponent}

\begin{modelcomponent}{Freeness Laws}
\textbf{No junk.}
\[
  t\in T_{\Sigma}(X)
  \Longrightarrow
  t\in X\ \vee\
  \exists f\in\Sigma\,\exists\bar t\,\bigl(t=f(\bar t)\bigr).
\]
\textbf{No confusion.}
\[
  x\neq f(\bar t),
  \qquad
  f(\bar s)=g(\bar t)\Longrightarrow f=g\wedge \bar s=\bar t.
\]
\end{modelcomponent}

\begin{modelconsequences}
  \modelconsequence{\textbf{Universal property.}
  \[
    \forall\mathcal A\;\forall\sigma:X\to\lvert\mathcal A\rvert\;
    \exists!\widehat\sigma:T_{\Sigma}(X)\to\lvert\mathcal A\rvert
  \]
  \[
    \operatorname{Hom}(\widehat\sigma,T_{\Sigma}(X),\mathcal A)
    \wedge
    \widehat\sigma\circ\eta=\sigma.
  \]}
  \modelconsequence{\[
    \operatorname{eval}(\mathcal A)(\sigma):=\widehat\sigma.
  \]}
\end{modelconsequences}

\begin{modelcomponent}{Instantiations}
\begin{center}
\begin{tabularx}{\linewidth}{@{}lXXXX@{}}
\toprule
use & \(\Sigma\) & \(X\) & target \(\mathcal A\) & seed \(\sigma\) \\
\midrule
formula eval &
\(\Omega\) &
atoms &
\(\mathbf 2\) &
atom truth values \\
parse &
\(\Omega\) &
\(X\) &
\(T_{\Omega}(X)\) &
\(\eta\) \\
term eval &
\(\Sigma_{\mathrm{term}}\) &
\(\mathsf{Var}\) &
\(\mathcal A_{\mathrm{term}}\) &
assignment \(s\) \\
\(\mathbb N\) preview &
\(\{0^{(0)},S^{(1)}\}\) &
\(\varnothing\) &
\((C,c_0,\operatorname{succ})\) &
none \\
\bottomrule
\end{tabularx}
\end{center}
\end{modelcomponent}

\begin{modelcomponent}{Worked Connective Instance}
\[
  \Omega=\{\neg^{(1)},\wedge^{(2)},\vee^{(2)},\to^{(2)},\leftrightarrow^{(2)}\},
\]
\[
  \mathbf 2=(\{0,1\},\neg^{\mathbf 2},\wedge^{\mathbf 2},
  \vee^{\mathbf 2},\to^{\mathbf 2},\leftrightarrow^{\mathbf 2}).
\]
\end{modelcomponent}

\begin{modelcomponent}{Boundary}
The machine fixes one seed. Connective and term evaluation are exact
instances. FOL quantifiers range over the assignment-indexed family
\((\sigma_s)_s\):
\[
  \mathcal A,s\models\exists x\,\varphi
  \Longleftrightarrow
  \exists a\in A,\;\mathcal A,s[x\mapsto a]\models\varphi.
\]
That is the downstream cylindric/polyadic layer. The absolutely free
\(T_{\Sigma}(X)\) is also distinct from the Lindenbaum--Tarski quotient by
\(\vdash\)-equivalence.
\end{modelcomponent}
\end{modelbox}

\begin{modelbox}{Consequence}
\modeldepends{Satisfaction relation \(\models\subseteq\mathsf{Mod}\times\mathsf{Form}\)}

\begin{modelcomponent}{Parameter}
A class of models \(\mathsf{Mod}\) and a satisfaction relation
\[
  m\models\varphi
  \qquad(m\in\mathsf{Mod}).
\]
\end{modelcomponent}

\begin{modelcomponent}{Semantic Consequence}
\[
  \Gamma\models\varphi
  \quad:\Longleftrightarrow\quad
  \forall m\in\mathsf{Mod}\,
  \bigl((\forall\gamma\in\Gamma,\;m\models\gamma)\to m\models\varphi\bigr).
\]
\[
  \mathrm{Cn}_{\models}(\Gamma):=\{\varphi:\Gamma\models\varphi\},
  \qquad
  \models\varphi:\Longleftrightarrow \varnothing\models\varphi.
\]
\[
  \mathrm{Th}(m):=\{\varphi:m\models\varphi\}.
\]
\end{modelcomponent}

\begin{modelcomponent}{Bindings}
\begin{center}
\begin{tabularx}{\linewidth}{@{}lXX@{}}
\toprule
logic & \(\mathsf{Mod}\) & \(m\models\varphi\) \\
\midrule
prop & valuations \(v\) & \(\widehat v(\varphi)=1\) \\
FOL open & pairs \((\mathcal A,s)\) & \(\mathcal A,s\models\varphi\) \\
FOL sentences & structures \(\mathcal A\) & \(\mathcal A\models\sigma\) \\
\bottomrule
\end{tabularx}
\end{center}
\end{modelcomponent}
\end{modelbox}

\section*{Propositional Logic}

\begin{modelbox}{Theory (Propositional Logic)}
\modeldepends{Propositional language; formula; Consequence; \(\mathrm{Cn}_{\vdash}\)}

\begin{modelcomponent}{Language}
\[
  \mathcal L_{\mathsf P}
  =
  (\mathsf P,\neg,\wedge,\vee,\to,\leftrightarrow),
  \qquad
  \mathsf{Form}_{\mathsf P}.
\]
\end{modelcomponent}

\begin{modelcomponent}{Presented Theories}
\[
  T:=\mathrm{Cn}_{\vdash}(\Sigma),
  \qquad
  \Sigma\subseteq\mathsf{Form}_{\mathsf P},
  \qquad
  \mathrm{CPL}:=\mathrm{Cn}_{\vdash}(\varnothing).
\]
\end{modelcomponent}

\begin{modelconsequences}
  \modelconsequence{\[
    \Gamma\vdash_{\mathrm{CPL}}\varphi
    \quad\Longleftrightarrow\quad
    \Gamma\models_{\mathrm{CPL}}\varphi
    \qquad(\text{completeness}).
  \]}
\end{modelconsequences}
\end{modelbox}

\begin{modelbox}{Model (Classical Propositional Logic)}
\begin{model}[Classical Propositional Logic]\label{model:classical-propositional-logic-sample}
\modeldepends{Propositional language; valuation; Free \(\Sigma\)-Algebra and Evaluation; Consequence}

\begin{modelcomponent}{Base Case of Evaluation}
The seed is the model datum itself.
\[
  v:\mathsf P\to\lvert\mathbf 2\rvert,
  \qquad
  \lvert\mathbf 2\rvert=\{0,1\}.
\]
\end{modelcomponent}

\begin{modelcomponent}{Binding}
\[
  \Sigma:=\Omega,
  \qquad
  X:=\mathsf P,
  \qquad
  \text{target}:=\mathbf 2,
  \qquad
  \sigma:=v,
  \qquad
  \widehat v:=\operatorname{eval}(\mathbf 2)(v).
\]
\end{modelcomponent}

\begin{modelbridge}
  \bridgeclause{\[
    v\models\varphi
    \quad:\Longleftrightarrow\quad
    \widehat v(\varphi)=1.
  \]}
\end{modelbridge}

\begin{modelconsequences}
  \modelconsequence{Connective recursion is inherited from
  \emph{Free \(\Sigma\)-Algebra and Evaluation}.}
  \modelconsequence{Semantic consequence, validity, and \(\mathrm{Th}(v)\) are
  inherited from \emph{Consequence}.}
\end{modelconsequences}
\end{model}
\end{modelbox}

\begin{modelbox}{Definitional Extension (Propositional Logic)}
\modeldepends{CPL; truth-function; conservative extension}

\begin{modelcomponent}{Admissibility}
A connective \(\star\) of arity \(n\) is admissible by a defining equivalence:
\[
  \star(\bar\varphi)\leftrightarrow\theta(\bar\varphi),
  \qquad
  \widehat v(\star(\bar\varphi))
  =
  F_{\star}(\widehat v(\varphi_1),\ldots,\widehat v(\varphi_n)).
\]
There is no \(\exists!\) obligation because formulas already evaluate to truth
values.
\end{modelcomponent}

\begin{modelcomponent}{Instances}
\[
  \top:\equiv(p_0\to p_0),
  \qquad
  \bot:\equiv\neg\top,
\]
\[
  \varphi\oplus\psi
  :\equiv
  ((\varphi\vee\psi)\wedge\neg(\varphi\wedge\psi)).
\]
\end{modelcomponent}

\begin{modelconsequences}
  \modelconsequence{\[
    \mathrm{CPL}^{+}\vdash\sigma
    \quad\Longleftrightarrow\quad
    \mathrm{CPL}\vdash\sigma
    \qquad(\sigma\in\mathcal L_{\mathsf P}).
  \]}
  \modelconsequence{Every extended formula has a base-language equivalent.}
\end{modelconsequences}
\end{modelbox}

\clearpage
\section*{Predicate Logic}

\begin{modelbox}{Theory (First-Order Logic)}
\modeldepends{Signature; term; formula; sentence; Consequence; \(\mathrm{Cn}_{\vdash}\)}

\begin{modelcomponent}{Language}
\[
  \mathcal L=(\mathsf{Const}_{\mathcal L},
  \mathsf{Func}_{\mathcal L},
  \mathsf{Rel}_{\mathcal L},=).
\]
\[
  \mathsf{Term}_{\mathcal L},
  \qquad
  \mathsf{Form}_{\mathcal L},
  \qquad
  \mathsf{Sent}_{\mathcal L}.
\]
\end{modelcomponent}

\begin{modelcomponent}{Presented Theory}
\[
  T:=\mathrm{Cn}_{\vdash}(\Sigma),
  \qquad
  \Sigma\subseteq\mathsf{Sent}_{\mathcal L}.
\]
\end{modelcomponent}

\begin{modelconsequences}
  \modelconsequence{\[
    \Gamma\vdash_{\mathrm{FOL}}\varphi
    \quad\Longleftrightarrow\quad
    \Gamma\models_{\mathrm{FOL}}\varphi
    \qquad(\text{completeness}).
  \]}
\end{modelconsequences}
\end{modelbox}

\begin{modelbox}{Model (First-Order Logic)}
\begin{model}[First-Order Logic]\label{model:first-order-logic-sample}
\modeldepends{\(\mathcal L\)-structure; interpretation; assignment; Free \(\Sigma\)-Algebra and Evaluation; Consequence}

\begin{modelcomponent}{Inductive Case of Evaluation}
FOL runs the shared machine twice: once for terms, then once for quantifier-free
formulas. Quantifiers sit above that fixed-seed machine.
\end{modelcomponent}

\begin{modelcomponent}{Structure and Interpretation}
\[
  \mathcal A=\langle A,I\rangle,
  \qquad
  A\neq\varnothing.
\]
\[
  I(c)\in A,
  \qquad
  I(f):A^n\to A,
  \qquad
  I(R)\subseteq A^n,
  \qquad
  =\text{ is logical equality on }A.
\]
\end{modelcomponent}

\begin{modelcomponent}{Two Evaluations, Stacked}
\[
  s:\mathsf{Var}\to A
  \xrightarrow{\operatorname{eval}(\mathcal A_{\mathrm{term}})}
  \llbracket\cdot\rrbracket_{\mathcal A,s}
  \xrightarrow{\text{relation/equality tests}}
  \sigma_s
  \xrightarrow{\operatorname{eval}(\mathbf 2)}
  \models_{\mathrm{qf}}.
\]
\end{modelcomponent}

\begin{modelcomponent}{[feeds seed] Term Evaluation}
\[
  \Sigma_{\mathrm{term}}
  :=
  \mathsf{Const}_{\mathcal L}\cup\mathsf{Func}_{\mathcal L},
  \qquad
  \mathcal A_{\mathrm{term}}
  :=
  (A,(I f)_{f\in\mathsf{Func}},(I c)_{c\in\mathsf{Const}}).
\]
\[
  \llbracket\cdot\rrbracket_{\mathcal A,s}
  :=
  \operatorname{eval}(\mathcal A_{\mathrm{term}})(s).
\]
\end{modelcomponent}

\begin{modelcomponent}{[seed, computed] Atomic Seed}
\[
  \sigma_s\bigl(R(\bar t)\bigr)
  :=
  \bigl[
    (\llbracket t_1\rrbracket_{\mathcal A,s},\ldots,
    \llbracket t_n\rrbracket_{\mathcal A,s})
    \in I(R)
  \bigr],
\]
\[
  \sigma_s(t_1=t_2)
  :=
  \bigl[
    \llbracket t_1\rrbracket_{\mathcal A,s}
    =
    \llbracket t_2\rrbracket_{\mathcal A,s}
  \bigr].
\]
\end{modelcomponent}

\begin{modelcomponent}{[connective eval, inherited] Quantifier-Free Truth}
\[
  \mathcal A,s\models\varphi
  \quad:\Longleftrightarrow\quad
  \operatorname{eval}(\mathbf 2)(\sigma_s)(\varphi)=1
  \qquad(\varphi\text{ quantifier-free}).
\]
\end{modelcomponent}

\begin{modelcomponent}{[\(\Omega^{+}\) extension] Quantifiers}
\[
  \mathcal A,s\models\exists x\,\varphi
  \quad\Longleftrightarrow\quad
  \exists a\in A,\;\mathcal A,s[x\mapsto a]\models\varphi.
\]
\[
  \mathcal A,s\models\forall x\,\varphi
  \quad\Longleftrightarrow\quad
  \forall a\in A,\;\mathcal A,s[x\mapsto a]\models\varphi.
\]
\end{modelcomponent}

\begin{modelcomponent}{Sentence Satisfaction}
For \(\sigma\in\mathsf{Sent}_{\mathcal L}\), \(\mathcal A,s\models\sigma\) is
independent of \(s\) by the coincidence lemma, so
\[
  \mathcal A\models\sigma
  :=
  \mathcal A,s\models\sigma
  \qquad(s\text{ arbitrary}).
\]
\end{modelcomponent}

\begin{modelconsequences}
  \modelconsequence{Semantic consequence, \(\mathcal A\models T\), and
  \(\mathrm{Th}(\mathcal A)\) are inherited from \emph{Consequence}.}
  \modelconsequence{Prop is the base case: seed \(=v\). FOL is the inductive
  case: \(s\) is evaluated into term values, then relation tests compute
  \(\sigma_s\).}
\end{modelconsequences}
\end{model}
\end{modelbox}

\begin{modelbox}{Definitional Extension (First-Order Logic)}
\modeldepends{Theory; expansion; reduct; conservativity; unique existence}

\begin{modelcomponent}{Logical Abbreviation}
\[
  \exists!y\,\varphi(y)
  \;:\equiv\;
  \exists y\bigl(\varphi(y)\wedge
  \forall y'(\varphi(y')\to y'=y)\bigr).
\]
\end{modelcomponent}

\begin{modelcomponent}{Admissibility}
\textbf{Relation symbol \(R\), arity \(n\):}
\[
  \forall\bar x\,(R(\bar x)\leftrightarrow\varphi(\bar x)).
\]
\textbf{Function symbol \(f\), arity \(n\):}
\[
  T\vdash \forall\bar x\,\exists!y\,\psi(\bar x,y),
  \qquad
  \forall\bar x\,\forall y\,(f(\bar x)=y\leftrightarrow\psi(\bar x,y)).
\]
\textbf{Constant symbol \(c\):}
\[
  T\vdash \exists!y\,\psi(y),
  \qquad
  \forall y\,(c=y\leftrightarrow\psi(y)).
\]
\end{modelcomponent}

\begin{modelconsequences}
  \modelconsequence{\[
    T^{+}\vdash\sigma
    \quad\Longleftrightarrow\quad
    T\vdash\sigma
    \qquad(\sigma\in\mathcal L).
  \]}
  \modelconsequence{\[
    \mathcal A\models T
    \Longrightarrow
    \exists!\mathcal A^{+}\,
    (\mathcal A^{+}\!\upharpoonright\mathcal L=\mathcal A).
  \]}
\end{modelconsequences}
\end{modelbox}

\end{document}
